% Test de bout en bout des modules syllabus (mod_syllabus, mod_syllabusstaff, mod_syllabusstudent).
% Compilation : latexmk -pdf test-syllabus.tex
\documentclass[11pt,a4paper]{article}
\usepackage[french]{babel}
\usepackage[T1]{fontenc}
\usepackage[utf8]{inputenc}
\usepackage{lmodern,microtype,booktabs,tabularx,xcolor,hyperref,array}
\usepackage[margin=2.2cm]{geometry}
\hypersetup{colorlinks=true, linkcolor=blue!50!black, urlcolor=blue!50!black}
\newcommand{\code}[1]{\texttt{\detokenize{#1}}}
\title{Test des modules syllabus et de leur évaluation}
\author{Sébastien Lefebvre}
\date{Novembre 2026 -- version 1.0.0 des trois modules}

\begin{document}
\maketitle

\section{Objet et méthode}

Les trois nouveaux modules -- \code{mod_syllabus} (fiche et évaluation de chaque module
d'enseignement), \code{mod_syllabusstaff} (suivi des enseignants et de la DEVE) et
\code{mod_syllabusstudent} (« Mes modules » des étudiants) -- ont été exercés par leurs pages
web, comme par des navigateurs (\code{docs/audit/banc/syllabus_interactions.py}), sur le banc
habituel (4 c\oe urs, nginx, PHP-FPM~8.3, PostgreSQL~16, Moodle~4.5). Chaque page est contrôlée
(erreur PHP, exception, code 500), puis la base.

\paragraph{Données} (\code{syllabus_setup.php}) : une catégorie de 12 cours de module, chacun
avec son activité « Syllabus du module » et deux responsables (24 enseignants) ; 200 étudiants
inscrits aux 12 cours ; la DEVE gestionnaire de chaque cours ; un cours d'équipe avec le suivi
des modules ; « Mes modules » dans le cours de la promotion.

\section{Scénario et contrôles}

\begin{center}
\begin{tabularx}{\linewidth}{@{}X l@{}}
\toprule
\textbf{Étape} & \textbf{Résultat} \\
\midrule
DEVE : référentiel (relecture obligatoire, seuil de 5 réponses, 2 types, 4 questions
communes dont une pour le type « Clinique »), type de chaque module dans ses réglages & conforme \\
12 responsables en parallèle : fiche rédigée (rubriques, 3 séances) et soumise, 2 questions & 2,5 s \\
Fiches à relire listées ; 11 publiées, 1 renvoyée puis corrigée et publiée & conforme \\
Fenêtre d'évaluation des 12 modules fixée en lot & conforme \\
Tâche d'ouverture : questions figées, étudiants prévenus & 2\,400 messages, 20 s \\
200 étudiants (40 en parallèle) : « Mes modules » (12 modules chacun), puis 9 évaluations sur 10
& 2\,172 réponses en 99 s \\
Seconde réponse d'un étudiant & refusée \\
Une réponse anonyme par participation, réponses obligatoires toutes présentes & 2\,172 / 2\,172 \\
Table des réponses : ni étudiant ni heure (\code{id}, \code{syllabusid}, \code{daysubmitted}) &
conforme \\
Questions par module : 3 communes, + 1 pour les modules « Clinique », + 2 de l'enseignant &
5 ou 6 \\
Résultats du module cachés au responsable avant la fermeture, visibles après ; export CSV &
conforme \\
Synthèse comparée de la DEVE (moyennes par question commune), liste filtrée, fiche PDF &
conforme \\
Anomalies dans les pages & 0 \\
\bottomrule
\end{tabularx}
\end{center}

Résultat : \textbf{23 contrôles sur 23}. S'y ajoutent 22 tests unitaires (droits, circuit de
relecture, figeage des questions, anonymat, seuil, visibilité des résultats, tâche,
réinitialisation, sauvegarde et restauration, RGPD, synthèses) et le test de conformité RGPD
de Moodle, réussi pour les trois modules.

\section{Défaut trouvé et corrigé}

Les captures du guide ont montré le formulaire d'évaluation avec une réponse \textbf{déjà
cochée} : le premier choix d'une question à choix (et « Non » pour une question oui/non). Le
bouton de valeur 0 était coché d'office, une valeur absente étant égale à 0 pour PHP. Un
étudiant pressé envoyait donc ce choix sans l'avoir fait, et le contrôle des réponses
obligatoires ne jouait plus pour ces questions. Les boutons portent désormais des valeurs
décalées d'une unité, ramenées à la réponse à l'enregistrement ; un test unitaire vérifie
qu'aucun bouton n'est coché et que les valeurs se ramènent aux réponses. Le scénario a été
rejoué après correction (23/23), et les réponses enregistrées gardent leurs valeurs (0 ou 1
pour oui/non, 0 à 2 pour un choix parmi trois).

Le scénario a aussi révélé une erreur dans le banc lui-même : les scripts appelaient un
programme de tâche inexistant au lieu de \code{admin/cli/scheduled_task.php}. Corrigé ; le scénario de liaison avec « Gestion des
stages », qui en dépendait pour sa tâche de nuit, a été rejoué (14/14).

\section{Limite connue}

L'anonymat vaut pour toutes les pages et tous les exports. En base, la participation et la
réponse sont écrites dans la même transaction : l'ordre des enregistrements permettrait, à qui
lit directement la base, un rapprochement approximatif. L'accès à la base doit donc rester
restreint ; le guide le signale.

\section{Reproduire}

{\small
\begin{verbatim}
cd docs/audit/banc
sudo -u www-data php syllabus_setup.php site-x.json > syllabus.json
python3 syllabus_interactions.py site-x.json syllabus.json > syllabus-resultats.json
sudo -u www-data php syllabus_setup.php site-x.json reset     # pour rejouer
\end{verbatim}
}

\end{document}
